\subsubsection*{2.7.}
态射 (2.6.1) 与 (2.6.2) 导出
\[
c:\varepsilon^*R\Psi
\longrightarrow R\Psi_{\mathrm{cl}}\circ\varepsilon^*
\]
和
\[
c_\eta:\varepsilon^*R\Psi_\eta
\longrightarrow R\Psi_{\eta\,\mathrm{cl}}\circ\varepsilon^*.
\]

\subsubsection*{定理 2.8.}
设 $K\in\operatorname{Ob}D(X,\Lambda)$（分别为 $K\in\operatorname{Ob}D(X-X_0,\Lambda)$）是 $\Lambda$-模层的复形，其上同调层 $\underline H^i(K)$ 可构造。则比较态射
\[
c:\varepsilon^*R\Psi(K)
\longrightarrow R\Psi_{\mathrm{cl}}(\varepsilon^*K)
\]
（分别为
\[
c:\varepsilon^*R\Psi_\eta(K)
\longrightarrow R\Psi_{\eta\,\mathrm{cl}}(\varepsilon^*K))
\]
是同构。

只需对 $R\Psi_\eta$ 证明此断言。此外，2.7 的两端都是有限上同调维数的三角函子，故只需处理 $K$ 是置于 $0$ 次的单个层 $F$ 的情形。证明与有限性定理 XIII, 2.3.1 的证明完全类似。由于 $R\Psi$ 与 $R\Psi_{\mathrm{cl}}$ 具有相同的变性性质，那里使用的逐步化简（允许缩小 $S$）将问题归结为 $X$ 正则、$f^{-1}(0)$ 是 $X$ 中的正规交叉除子且 $F=\mathbb Z/\ell\mathbb Z$ 的情形。此时比较第 I 讲 3.3 与其超越对应结论即可。

